\appendix
\section*{Appendix D — Spectral Ecology}
\addcontentsline{toc}{section}{Appendix D — Spectral Ecology}

\section{Spectral Framework}

\subsection{Laplacian Definition}
The spectral ecology layer is built on the discrete Laplacian
\[
L = D - W,
\]
where $D$ is the degree matrix and $W$ is the curvature-aware adjacency
matrix. The Laplacian is symmetric and positive semi-definite.

\subsection{Eigenstructure}
The Laplacian admits eigenpairs
\[
L u_k = \lambda_k u_k,
\]
where $u_k$ is the $k$-th eigenvector and $\lambda_k$ the corresponding
eigenvalue.

\subsection{Ecological Interpretation}
Eigenvectors represent monster families:
\[
\text{MonsterFamily}_k = u_k,
\]
and eigenvalues represent ecological energy levels:
\[
\text{EnergyLevel}_k = \lambda_k.
\]

\section{Curvature-Induced Spectral Drift}

\subsection{Curvature Modification}
Holonomy curvature modifies adjacency weights:
\[
w_{ij}' = w_{ij} + H_{ij},
\]
producing a warped Laplacian
\[
L' = D' - W'.
\]

\subsection{Eigenvalue Drift}
Perturbation of the Laplacian produces spectral drift:
\[
\lambda_k' = \lambda_k + \delta \lambda_k.
\]

\subsection{Eigenvector Drift}
Eigenvectors shift under curvature:
\[
u_k' = u_k + \delta u_k.
\]

\section{Definitions}

\begin{definition}[Spectral Ecology]
Spectral ecology is the study of how eigenvectors and eigenvalues of the
Laplacian encode monster families and ecological energy levels.
\end{definition}

\begin{definition}[Spectral Drift]
Spectral drift is the change in eigenstructure induced by curvature:
\[
(\lambda_k, u_k) \mapsto (\lambda_k', u_k').
\]
\end{definition}

\begin{definition}[Ecological Stability]
A monster family is stable if
\[
\|u_k' - u_k\| \le \epsilon
\]
for small curvature perturbations $\epsilon$.
\end{definition}

\begin{definition}[Energy Level]
The ecological energy level of a monster family is given by its
eigenvalue $\lambda_k$.
\end{definition}

\section{Theorems}

\begin{theorem}[Spectral Stability Under Small Perturbations]
Let $L' = L + \epsilon K$ for small $\epsilon$. Then
\[
|\lambda_k' - \lambda_k| \le \epsilon \|K\|.
\]
\end{theorem}

\begin{proof}
By classical perturbation theory for symmetric matrices,
\[
\lambda_k' = \lambda_k + u_k^\top (\epsilon K) u_k + O(\epsilon^2),
\]
so the deviation is bounded by $\epsilon \|K\|$.
\end{proof}

\begin{theorem}[Eigenvector Stability]
For small curvature perturbations,
\[
\|u_k' - u_k\| \le \epsilon C,
\]
for some constant $C$ depending on spectral gaps.
\end{theorem}

\begin{proof}
Eigenvector perturbation bounds follow from Davis–Kahan theory:
\[
\|u_k' - u_k\|
    \le \frac{\epsilon \|K\|}{\min_{j \neq k} |\lambda_k - \lambda_j|}.
\]
\end{proof}

\begin{theorem}[Curvature Causes Ecological Evolution]
If curvature is non-zero, then at least one eigenpair changes:
\[
(\lambda_k', u_k') \neq (\lambda_k, u_k).
\]
\end{theorem}

\begin{proof}
Curvature modifies adjacency weights, so $L' \neq L$. Since the Laplacian
is symmetric, any non-zero perturbation changes at least one eigenpair.
\end{proof}

\section{Spectral–Laplacian Interaction}

\subsection{Warped Laplacian}
Curvature produces a warped Laplacian
\[
L' = D' - W',
\]
which shifts spectral ecology.

\subsection{Ecological Consequences}
Spectral drift produces:
\begin{itemize}
    \item monster-family evolution,
    \item faction-biased ecological zones,
    \item narrative-tension propagation,
    \item curvature-sensitive rare species.
\end{itemize}

\section{Discussion}
This appendix formalizes the spectral ecology layer, including
eigenstructure interpretation, curvature-induced drift, stability
analysis, and ecological consequences. These structures integrate with
holonomy curvature, Laplacian geometry, and reversible quest dynamics in
the main preprint.
